\chapter{Grundlagen der Kernphysik}

\section{Ziel dieses Kapitels}

Dieses Kapitel führt die Grundbegriffe der Kernphysik ein.
Es behandelt Nuklide, Isotope, Isobare, Isotone, die Nuklidkarte, Kernradien,
Massendefekt, Bindungsenergie und die erste qualitative Idee von Kernstabilität.

\begin{notebox}
Dieses Kapitel ist die Grundlage für spätere Themen wie:
\begin{itemize}
    \item Rutherfordstreuung und Formfaktoren,
    \item Weizsäcker-Massenformel,
    \item Massental,
    \item Kernmodelle,
    \item radioaktive Zerfälle,
    \item Kernspaltung und Kernfusion.
\end{itemize}
\end{notebox}

\section{Atomkerne, Protonen und Neutronen}

Ein Atomkern besteht aus Protonen und Neutronen.
Protonen und Neutronen werden zusammen als \emph{Nukleonen} bezeichnet.

\begin{definitionbox}
Die \emph{Protonenzahl} $Z$ gibt an, wie viele Protonen ein Kern enthält.

Die \emph{Neutronenzahl} $N$ gibt an, wie viele Neutronen ein Kern enthält.

Die \emph{Massenzahl} $A$ gibt die Gesamtzahl der Nukleonen an:
\[
    A = Z + N.
\]
\end{definitionbox}

\begin{conceptbox}
Die Protonenzahl $Z$ bestimmt das chemische Element.

Beispiele:
\[
    Z=1 \Rightarrow \text{Wasserstoff},
    \qquad
    Z=6 \Rightarrow \text{Kohlenstoff},
    \qquad
    Z=92 \Rightarrow \text{Uran}.
\]

Die Neutronenzahl $N$ bestimmt dagegen, welche Kernsorte dieses Elements vorliegt.
Unterschiedliche Neutronenzahlen können zu unterschiedlicher Stabilität führen.
\end{conceptbox}

\begin{formulabox}
\[
    A = Z + N
\]
und damit
\[
    N = A - Z.
\]
\end{formulabox}

\begin{examtipbox}
Bei jedem Nuklid zuerst bestimmen:
\[
    Z,\qquad A,\qquad N=A-Z.
\]
Die Massenzahl $A$ ist nicht die Neutronenzahl.
\end{examtipbox}

\section{Schreibweise von Nukliden}

Ein Nuklid wird meistens in der Form
\[
    {}^A_Z X
\]
geschrieben. Dabei ist $X$ das Elementsymbol.

\begin{definitionbox}
Ein \emph{Nuklid} ist eine Kernsorte mit eindeutig festgelegter Protonenzahl $Z$
und Neutronenzahl $N$.

Da
\[
    A=Z+N
\]
gilt, kann man ein Nuklid auch durch $Z$ und $A$ eindeutig festlegen.
\end{definitionbox}

\begin{examplebox}
Betrachte das Nuklid
\[
    {}^{238}_{92}\mathrm{U}.
\]
Dann ist
\[
    Z=92,
    \qquad
    A=238.
\]
Die Neutronenzahl ist daher
\[
    N=A-Z=238-92=146.
\]
Der Kern enthält also $92$ Protonen und $146$ Neutronen.
\end{examplebox}

\section{Nuklidsorten}

Die wichtigsten Nuklidsorten sind Isotope, Isobare, Isotone und Kernisomere.
Diese Begriffe beschreiben, welche der Zahlen $Z$, $N$ und $A$ gleich bleiben.

\subsection{Isotope}

\begin{definitionbox}
\emph{Isotope} sind Nuklide mit gleicher Protonenzahl $Z$, aber unterschiedlicher
Neutronenzahl $N$.

Sie gehören also zum gleichen chemischen Element, haben aber unterschiedliche
Massenzahlen $A$.
\end{definitionbox}

\begin{conceptbox}
Isotope sind chemisch sehr ähnlich, weil das chemische Verhalten vor allem durch
die Elektronenhülle bestimmt wird.
Bei neutralen Atomen ist die Elektronenzahl gleich der Protonenzahl $Z$.

Kernphysikalisch können Isotope aber sehr unterschiedlich sein.
Ein Isotop kann stabil sein, ein anderes radioaktiv.
\end{conceptbox}

\begin{examplebox}
Die Kohlenstoffisotope
\[
    {}^{12}_{6}\mathrm{C},
    \qquad
    {}^{13}_{6}\mathrm{C},
    \qquad
    {}^{14}_{6}\mathrm{C}
\]
haben alle
\[
    Z=6.
\]
Die Neutronenzahlen sind:
\[
    N({}^{12}\mathrm{C})=12-6=6,
\]
\[
    N({}^{13}\mathrm{C})=13-6=7,
\]
\[
    N({}^{14}\mathrm{C})=14-6=8.
\]
Sie sind also Isotope desselben Elements.
\end{examplebox}

\subsection{Isobare}

\begin{definitionbox}
\emph{Isobare} sind Nuklide mit gleicher Massenzahl $A$, aber unterschiedlicher
Protonenzahl $Z$.
\end{definitionbox}

\begin{conceptbox}
Isobare haben gleich viele Nukleonen, aber eine unterschiedliche Aufteilung in
Protonen und Neutronen.

Dadurch ändern sich:
\begin{itemize}
    \item die Coulomb-Abstoßung zwischen den Protonen,
    \item das Verhältnis von Neutronen zu Protonen,
    \item die Kernmasse,
    \item die Stabilität.
\end{itemize}
\end{conceptbox}

\begin{examplebox}
Die Kerne
\[
    {}^{14}_{6}\mathrm{C}
    \qquad\text{und}\qquad
    {}^{14}_{7}\mathrm{N}
\]
haben beide
\[
    A=14.
\]
Sie besitzen aber verschiedene Protonenzahlen:
\[
    Z({}^{14}\mathrm{C})=6,
    \qquad
    Z({}^{14}\mathrm{N})=7.
\]
Damit sind sie Isobare.
\end{examplebox}

\subsection{Isotone}

\begin{definitionbox}
\emph{Isotone} sind Nuklide mit gleicher Neutronenzahl $N$, aber unterschiedlicher
Protonenzahl $Z$.
\end{definitionbox}

\begin{examplebox}
Betrachte
\[
    {}^{14}_{6}\mathrm{C}
    \qquad\text{und}\qquad
    {}^{15}_{7}\mathrm{N}.
\]
Dann gilt:
\[
    N({}^{14}\mathrm{C})=14-6=8,
\]
\[
    N({}^{15}\mathrm{N})=15-7=8.
\]
Beide Kerne haben also dieselbe Neutronenzahl.
Sie sind Isotone.
\end{examplebox}

\subsection{Kernisomere}

\begin{definitionbox}
Ein \emph{Kernisomer} ist ein langlebiger angeregter Zustand eines Atomkerns.

Ein Kernisomer besitzt dieselbe Protonenzahl $Z$ und dieselbe Massenzahl $A$
wie der zugehörige Grundzustand, aber eine andere innere Energie.
\end{definitionbox}

\begin{conceptbox}
Ein Kernisomer ist kein anderes Element und kein anderes Isotop.
Es ist derselbe Kern in einem metastabilen angeregten Zustand.

Solche Zustände können zum Beispiel durch Aussenden von Gammastrahlung zerfallen.
\end{conceptbox}

\subsection{Vergleich der Nuklidsorten}

\begin{center}
\begin{tabular}{|l|l|l|p{6.4cm}|}
\hline
Begriff & Gleich & Verschieden & Bedeutung \\
\hline
Isotope & $Z$ & $N$, $A$ & gleiches chemisches Element \\
\hline
Isobare & $A$ & $Z$, $N$ & gleiche Gesamtzahl an Nukleonen \\
\hline
Isotone & $N$ & $Z$, $A$ & gleiche Neutronenzahl \\
\hline
Isomere & $Z$, $A$ & Energiezustand & gleicher Kern in angeregtem Zustand \\
\hline
\end{tabular}
\end{center}

\begin{examtipbox}
Merke:
\[
    \text{Isotope: } Z \text{ gleich},
\]
\[
    \text{Isobare: } A \text{ gleich},
\]
\[
    \text{Isotone: } N \text{ gleich}.
\]
\end{examtipbox}

\section{Nuklidkarte}

Die Nuklidkarte ordnet Kerne nach Protonenzahl und Neutronenzahl.
Sie ist eine Art Landkarte der Kernstabilität.

\begin{definitionbox}
Eine \emph{Nuklidkarte} ist eine Darstellung von Nukliden in einem Koordinatensystem,
in dem typischerweise die Protonenzahl $Z$ und die Neutronenzahl $N$ aufgetragen werden.
\end{definitionbox}

\begin{conceptbox}
Stabile Kerne liegen nicht zufällig in der Nuklidkarte.
Sie bilden ein charakteristisches Stabilitätsband.

Für leichte stabile Kerne gilt ungefähr:
\[
    N \approx Z.
\]
Für schwere stabile Kerne gilt meistens:
\[
    N > Z.
\]
\end{conceptbox}

\begin{conceptbox}
Der Grund für $N>Z$ bei schweren stabilen Kernen ist die Coulomb-Abstoßung zwischen Protonen.

Protonen stoßen sich elektrisch ab.
Neutronen tragen zur starken Kernbindung bei, erzeugen aber keine zusätzliche elektrische
Abstoßung.

Schwere Kerne benötigen daher zusätzliche Neutronen, um stabil zu bleiben.
\end{conceptbox}

\begin{notebox}
In der Nuklidkarte werden auch Zerfallstypen markiert.
Typische Zerfallstypen sind:
\[
    \beta^+,
    \qquad
    \varepsilon\text{-Einfang},
    \qquad
    \beta^-,
    \qquad
    \alpha,
    \qquad
    \text{spontane Spaltung}.
\]
Die genaue Behandlung dieser Zerfälle folgt im Kapitel über Radioaktivität.
\end{notebox}

\begin{examtipbox}
Wenn ein Kern zu viele Neutronen besitzt, ist häufig $\beta^-$-Zerfall relevant.

Wenn ein Kern zu viele Protonen besitzt, sind häufig $\beta^+$-Zerfall oder
Elektroneneinfang relevant.

Das ist eine qualitative Orientierung. Die genaue Entscheidung hängt von Massen
und Energiebilanzen ab.
\end{examtipbox}

\section{Kernradius}

Die typische Längenskala der Kernphysik ist der Femtometer:
\[
    1\,\si{fm}=10^{-15}\,\si{m}.
\]

\begin{definitionbox}
Der \emph{Kernradius} beschreibt die räumliche Ausdehnung eines Atomkerns.

Eine wichtige empirische Näherung ist:
\[
    R = r_0 A^{1/3},
\]
wobei
\[
    r_0 \approx \SI{1.2}{fm}
\]
ist.
\end{definitionbox}

\begin{formulabox}
\[
    R = r_0 A^{1/3},
    \qquad
    r_0 \approx \SI{1.2}{fm}.
\]
\end{formulabox}

\begin{derivationbox}
Nimmt man den Kern näherungsweise als Kugel an, gilt:
\[
    V = \frac{4\pi}{3}R^3.
\]
Da Kernmaterie näherungsweise konstante Dichte besitzt, ist das Kernvolumen
proportional zur Nukleonenzahl:
\[
    V \propto A.
\]
Damit folgt:
\[
    R^3 \propto A.
\]
Also:
\[
    R \propto A^{1/3}.
\]
Daraus ergibt sich die Näherung:
\[
    R = r_0 A^{1/3}.
\]
\end{derivationbox}

\begin{examplebox}
Schätze den Radius von ${}^{208}\mathrm{Pb}$.

Für Blei-208 gilt:
\[
    A=208.
\]
Damit:
\[
    R \approx \SI{1.2}{fm}\cdot 208^{1/3}.
\]
Mit
\[
    208^{1/3}\approx 5.92
\]
folgt:
\[
    R \approx \SI{7.1}{fm}.
\]
\end{examplebox}

\begin{examtipbox}
Wenn $A$ um den Faktor $8$ größer wird, wächst der Radius nur um
\[
    8^{1/3}=2.
\]
Der Radius verdoppelt sich also.
\end{examtipbox}

\section{Elektromagnetischer Radius}

Der Begriff Kernradius ist nicht eindeutig.
Je nach Messmethode erhält man unterschiedliche Radiusbegriffe.

\begin{definitionbox}
Der \emph{elektromagnetische Radius} beschreibt die räumliche Ausdehnung der
elektrischen Ladungsverteilung im Kern.
\end{definitionbox}

\begin{conceptbox}
Atomkerne besitzen keine harte Oberfläche.
Die Ladungsverteilung fällt am Kernrand über eine Randzone ab.

Wichtige Größen zur Beschreibung der Ladungsverteilung sind zum Beispiel:
\[
    R_{1/2}
\]
als Radius, bei dem die Dichte auf die Hälfte der zentralen Dichte gefallen ist,
und
\[
    \sqrt{\langle r^2\rangle}
\]
als mittlerer quadratischer Radius.
\end{conceptbox}

\begin{notebox}
Die genaue Beziehung zwischen Ladungsverteilung, elektromagnetischem Radius und
Formfaktor wird im Kapitel über Streuung und Formfaktoren behandelt.
\end{notebox}

\section{Rutherfordradius}

\begin{definitionbox}
Der \emph{Rutherfordradius} ist ein Radiusbegriff, der aus der Coulomb-Streuung
geladener Teilchen an Kernen motiviert ist.

Er hängt mit der kleinsten Annäherung eines geladenen Projektils an einen geladenen
Kern zusammen.
\end{definitionbox}

\begin{conceptbox}
Bei der Rutherfordstreuung wird die Ablenkung eines geladenen Projektils durch das
Coulomb-Potential des Kerns beschrieben.

Solange die gemessene Streuung der Rutherfordform folgt, sieht das Projektil im
Wesentlichen nur die Coulomb-Wechselwirkung.

Abweichungen von der Rutherfordstreuung zeigen, dass die endliche Kernausdehnung
oder die starke Wechselwirkung relevant wird.
\end{conceptbox}

\begin{notebox}
Der Rutherfordradius wird im Streuungskapitel genauer behandelt.
Hier ist nur wichtig, dass unterschiedliche Experimente unterschiedliche
Radiusbegriffe liefern können.
\end{notebox}

\section{Massen in der Kernphysik}

In der Kernphysik sind Massen und Energien direkt miteinander verbunden:
\[
    E = mc^2.
\]

\begin{definitionbox}
Die atomare Masseneinheit wird mit $\mathrm{u}$ bezeichnet.
Sie ist definiert als ein Zwölftel der Masse eines neutralen Kohlenstoff-12-Atoms:
\[
    1\,\mathrm{u} = \frac{1}{12}m({}^{12}\mathrm{C}).
\]
Numerisch gilt:
\[
    1\,\mathrm{u}c^2 \approx \SI{931.5}{MeV}.
\]
\end{definitionbox}

\begin{conceptbox}
Teilchen- und Kernmassen werden häufig als Energieäquivalente angegeben.
Statt eine Masse in Kilogramm anzugeben, schreibt man zum Beispiel:
\[
    m c^2 = \SI{931.5}{MeV}.
\]
Oder man gibt die Masse als
\[
    m = \SI{931.5}{MeV}/c^2
\]
an.
\end{conceptbox}

\begin{formulabox}
\[
    E = mc^2,
    \qquad
    1\,\mathrm{u}c^2 \approx \SI{931.5}{MeV}.
\]
\end{formulabox}

\section{Massendefekt}

Ein gebundener Atomkern ist leichter als die Summe seiner freien Nukleonen.

\begin{definitionbox}
Der \emph{Massendefekt} $\Delta m$ ist die Differenz zwischen der Summe der
Massen der freien Nukleonen und der Masse des gebundenen Kerns:
\[
    \Delta m = Zm_p + Nm_n - m_{\mathrm{Kern}}.
\]
\end{definitionbox}

\begin{formulabox}
\[
    \Delta m = Zm_p + Nm_n - m_{\mathrm{Kern}}.
\]
Dabei ist:
\begin{itemize}
    \item $m_p$ die Masse eines freien Protons,
    \item $m_n$ die Masse eines freien Neutrons,
    \item $m_{\mathrm{Kern}}$ die Masse des gebundenen Kerns.
\end{itemize}
\end{formulabox}

\begin{conceptbox}
Der Massendefekt bedeutet nicht, dass Masse verschwunden ist.
Die fehlende Masse entspricht der Bindungsenergie des Kerns.

Ein gebundenes System hat eine niedrigere Energie als die getrennten Bestandteile.
Nach
\[
    E=mc^2
\]
entspricht diese niedrigere Energie einer kleineren Masse.
\end{conceptbox}

\section{Bindungsenergie}

\begin{definitionbox}
Die \emph{Bindungsenergie} $B$ ist die Energie, die benötigt wird, um einen Kern
vollständig in freie Protonen und Neutronen zu zerlegen.

Gleichzeitig ist sie die Energie, die frei wird, wenn der Kern aus freien Nukleonen
gebildet wird.
\end{definitionbox}

\begin{formulabox}
\[
    B = \Delta m c^2.
\]
Mit dem Massendefekt:
\[
    B =
    \left(
        Zm_p + Nm_n - m_{\mathrm{Kern}}
    \right)c^2.
\]
\end{formulabox}

\begin{definitionbox}
Die \emph{Bindungsenergie pro Nukleon} ist
\[
    \frac{B}{A}.
\]
Sie gibt an, wie stark ein einzelnes Nukleon im Kern im Mittel gebunden ist.
\end{definitionbox}

\begin{conceptbox}
Die gesamte Bindungsenergie $B$ wächst grob mit der Größe des Kerns.
Ein großer Kern hat daher schon deshalb eine große Gesamtbindungsenergie, weil er
viele Nukleonen enthält.

Zum Vergleich verschiedener Kerne ist deshalb die Bindungsenergie pro Nukleon
wichtiger:
\[
    \frac{B}{A}.
\]
\end{conceptbox}

\begin{examplebox}
Angenommen, ein Kern hat den Massendefekt
\[
    \Delta m = 0.0304\,\mathrm{u}.
\]
Dann ist:
\[
    B = \Delta m c^2.
\]
Mit
\[
    1\,\mathrm{u}c^2 \approx \SI{931.5}{MeV}
\]
folgt:
\[
    B = 0.0304 \cdot \SI{931.5}{MeV}
      \approx \SI{28.3}{MeV}.
\]
Falls
\[
    A=4
\]
ist, erhält man:
\[
    \frac{B}{A}
    =
    \frac{\SI{28.3}{MeV}}{4}
    \approx \SI{7.1}{MeV}.
\]
\end{examplebox}

\section{Bindungsenergie pro Nukleon und Energiegewinn}

\begin{conceptbox}
Die Kurve der Bindungsenergie pro Nukleon erklärt qualitativ, warum sowohl
Kernfusion als auch Kernspaltung Energie freisetzen können.

Bei leichten Kernen steigt $B/A$ mit zunehmender Massenzahl stark an.
Fusion leichter Kerne kann daher Energie freisetzen.

Bei sehr schweren Kernen fällt $B/A$ langsam wieder ab.
Spaltung schwerer Kerne in mittelschwere Kerne kann daher ebenfalls Energie freisetzen.
\end{conceptbox}

\begin{formulabox}
Für eine Kernreaktion ist der $Q$-Wert:
\[
    Q =
    \left(
        m_{\mathrm{Anfang}} - m_{\mathrm{Ende}}
    \right)c^2.
\]
Wenn
\[
    Q>0,
\]
wird Energie frei.

Wenn
\[
    Q<0,
\]
muss Energie zugeführt werden.
\end{formulabox}

\section{Isobaren, Massen und Stabilität}

In den Vorlesungsfolien werden Kerne mit gleicher Massenzahl $A$ betrachtet.
Das sind Isobare.

\begin{conceptbox}
Bei einer Isobarenreihe ist $A$ fest.
Wenn sich $Z$ ändert, ändert sich automatisch auch
\[
    N=A-Z.
\]

Die Kernmasse hängt davon ab, wie viele Protonen und Neutronen vorhanden sind.
Der stabile Kern einer Isobarenreihe liegt typischerweise beim Massenminimum.
\end{conceptbox}

\begin{examplebox}
Ein Beispiel für eine Isobarenreihe ist:
\[
    A=81.
\]
Dazu gehören zum Beispiel:
\[
    {}^{81}_{33}\mathrm{As},
    \quad
    {}^{81}_{34}\mathrm{Se},
    \quad
    {}^{81}_{35}\mathrm{Br},
    \quad
    {}^{81}_{36}\mathrm{Kr},
    \quad
    {}^{81}_{37}\mathrm{Rb}.
\]
Alle besitzen dieselbe Massenzahl $A=81$, aber verschiedene Protonenzahlen.
\end{examplebox}

\begin{notebox}
Die genaue Beschreibung, warum die Masse entlang einer Isobarenreihe ein Minimum
besitzt, erfolgt später mit der Weizsäcker-Massenformel.
Dabei spielen insbesondere Coulombterm und Asymmetrieterm eine wichtige Rolle.
\end{notebox}

\section{Qualitative Kernstabilität}

\begin{conceptbox}
Ob ein Kern stabil ist, hängt vom Zusammenspiel mehrerer Effekte ab:
\begin{itemize}
    \item starke Wechselwirkung zwischen Nukleonen,
    \item Coulomb-Abstoßung zwischen Protonen,
    \item Verhältnis von Neutronenzahl $N$ zu Protonenzahl $Z$,
    \item Schaleneffekte,
    \item Paarungseffekte.
\end{itemize}
\end{conceptbox}

\begin{notebox}
Diese Stabilitätsfragen werden in späteren Kapiteln genauer behandelt:
\begin{itemize}
    \item Weizsäcker-Massenformel,
    \item Massental,
    \item magische Zahlen,
    \item Kernmodelle,
    \item radioaktive Zerfälle.
\end{itemize}
\end{notebox}

\section{Mini-Zusammenfassung}

\begin{notebox}
Die wichtigsten Punkte dieses Kapitels sind:
\begin{itemize}
    \item Ein Nuklid ist durch $Z$ und $N$ eindeutig bestimmt.
    \item Es gilt
    \[
        A=Z+N.
    \]
    \item Isotope haben gleiches $Z$.
    \item Isobare haben gleiches $A$.
    \item Isotone haben gleiches $N$.
    \item Die Nuklidkarte zeigt stabile und instabile Kerne.
    \item Für schwere stabile Kerne gilt meistens $N>Z$.
    \item Der Kernradius skaliert näherungsweise wie
    \[
        R=r_0A^{1/3}.
    \]
    \item Der Massendefekt ist
    \[
        \Delta m = Zm_p + Nm_n - m_{\mathrm{Kern}}.
    \]
    \item Die Bindungsenergie ist
    \[
        B=\Delta mc^2.
    \]
    \item Die Bindungsenergie pro Nukleon $B/A$ ist wichtig zum Vergleich verschiedener Kerne.
\end{itemize}
\end{notebox}

\section{Prüfungscheck}

\begin{examtipbox}
Nach diesem Kapitel solltest du folgende Fragen beantworten können:
\begin{enumerate}
    \item Was bedeuten $Z$, $N$ und $A$?
    \item Wie berechnet man aus ${}^A_ZX$ die Neutronenzahl?
    \item Was ist ein Nuklid?
    \item Was ist der Unterschied zwischen Isotopen, Isobaren und Isotonen?
    \item Was ist ein Kernisomer?
    \item Was zeigt die Nuklidkarte?
    \item Warum besitzen schwere stabile Kerne meistens mehr Neutronen als Protonen?
    \item Warum gilt näherungsweise $R=r_0A^{1/3}$?
    \item Was ist der elektromagnetische Radius?
    \item Was ist der Rutherfordradius?
    \item Was ist der Massendefekt?
    \item Wie hängt die Bindungsenergie mit dem Massendefekt zusammen?
    \item Warum betrachtet man die Bindungsenergie pro Nukleon?
    \item Was bedeutet ein positiver $Q$-Wert?
    \item Warum ist bei Isobaren das Massenminimum wichtig für Stabilität?
\end{enumerate}
\end{examtipbox}